\subsection{Hermitian forms}

We recall the following definition given in Algebra (A, IX, § 3, No.~1):

DEFINITION 1. --- Let E be a vector space over the field K. A hermitian form (on the left) on E is a map f from E × E into K satisfying the following conditions (for \(x_{1}\) , \(x_{2}\) , x, \(y_{1}\) , \(y_{2}\) , y in E and \(\lambda\) , \(\mu\) in K):

\begin{enumerate}
\def\labelenumi{(\arabic{enumi})}
\tightlist
\item
\item
\end{enumerate}

\[
\begin{array}{c} \left\{ \begin{array}{l} f (x _ {1} + x _ {2}, y) = f (x _ {1}, y) + f (x _ {2}, y) \\ f (x, y _ {1} + y _ {2}) = f (x, y _ {1}) + f (x, y _ {2}) \end{array} \right. \\ \left\{ \begin{array}{l} f (\lambda x, y) = \overline {{\lambda}} f (x, y) \\ f (x, \mu y) = \mu f (x, y) \\ f (x, y) = \overline {{f (y , x)}}  . \end{array} \right. \end{array}\tag{3}
\]

When the field K is R, the notion of hermitian form on E reduces to that of symmetric bilinear form on \(E \times E\) (A, III, § 6, No.~3).

We note that the second condition (1) and the second condition (2) follow from the other three.

From (1) and (2) we deduce immediately that

\[
f \left(\sum_ {j} \lambda_ {j} x _ {j}, \sum_ {k} \mu_ {k} y _ {k}\right) = \sum_ {j, k} \bar {\lambda} _ {j} \mu_ {k} f \left(x _ {j}, y _ {k}\right).\tag{4}
\]

In particular, if \(\mathbf{E}\) is finite dimensional, and if \((e_j)_{1 \leqslant j \leqslant n}\) is a basis of \(\mathbf{E}\) , then for \(x = \sum_{j=1}^{n} \xi_j e_j\) and \(y = \sum_{j=1}^{n} \eta_j e_j\) , we have,

\[
f (x, y) = \sum_ {j, k} \alpha_ {j k} \overline {{{{\xi}}}} _ {j} \eta_ {k}
\]

with the notation \(\alpha_{jk} = f(e_j, e_k)\) ; moreover, relation (3) amounts to \(\alpha_{jk} = \overline{\alpha_{kj}}\) for every pair of indices \(j, k\) ; this implies in particular that the numbers \(\alpha_{jj}\) are real.

From (3), the number \(Q(x) = f(x, x)\) is real for all \(x \in E\) . Moreover, we immediately establish the following formulas, known as polarization formulas

\begin{enumerate}
\def\labelenumi{(\arabic{enumi})}
\setcounter{enumi}{4}
\tightlist
\item
\end{enumerate}

\[
4 f (x, y) = \sum_ {\varepsilon^ {2} = 1} \varepsilon \mathrm{Q} (x + \varepsilon y) \quad \text { if } \quad \mathbf {K} = \mathbf {R},\tag{6}
\]

\[
4 f (x, y) = \sum_ {\varepsilon^ {4} = 1, \varepsilon \in \mathbf {C}} \varepsilon \mathrm{Q} (x + \overline {{{\varepsilon}}} y) \quad \text { if } \quad \mathbf {K} = \mathbf {C}.
\]

Remark. --- We observe that formula (6) is valid for every sesquilinear form on \(E \times E\) (that is to say, for every function f satisfying (1) and (2), but not necessarily (3)). This remark shows that, when K = C, a sesquilinear form f such that \(f(x, x)\) is real for all \(x \in E\) is necessarily hermitian: relation (6) then gives \(\overline{f(y, x)} = f(x, y)\) since we have \(y + \varepsilon x = \varepsilon(x + \overline{\varepsilon}y)\) and \(\mathbf{Q}(\varepsilon z) = \mathbf{Q}(z)\) whenever \(\varepsilon^{4} = 1\) .

From the polarisation formulas, we have in particular,

PROPOSITION 1. --- If \(f\) is a hermitian form on \(\mathbf{E}\) , and \(\mathbf{M}\) a vector subspace of \(\mathbf{E}\) such that \(f(x, x) = 0\) for all \(x \in \mathbf{M}\) , then we also have \(f(x, y) = 0\) for every pair of points \(x, y\) in \(\mathbf{M}\) .

Let \(f\) be a hermitian form on \(\mathbf{E}\) ; the set \(\mathbf{N}\) of all \(x \in \mathbf{E}\) such that \(f(x, y) = 0\) for all \(y \in \mathbf{E}\) is a vector subspace of \(\mathbf{E}\) . It follows from (3) that, if \(x_1 \equiv x_2 \pmod{\mathbf{N}}\) and \(y_1 \equiv y_2 \pmod{\mathbf{N}}\) , we have \(f(x_1, y_1) = f(x_2, y_2)\) ; hence, on the quotient space \(\mathrm{E/N}\) we define a sesquilinear form \(f\) by putting \(\dot{f}(\dot{x}, \dot{y}) = f(x, y)\) for all \(x \in \dot{x}\) and all \(y \in \dot{y}\) ; it is clear that \(\dot{f}\) is hermitian and that the relation \(\langle \dot{f}(\dot{x}, \dot{y}) = 0 \text{ for all } \dot{y} \in \mathrm{E/N} \rangle\) implies \(\dot{x} = 0 \text{ in } \mathrm{E/N}\) , in other words (A, IX) \(\dot{f}\) is separating. We say that \(f\) is the separating hermitian form associated with \(f\) .

